\chapter{Debugging}

\usetikzlibrary{shapes.geometric}
\tikzset{pbox/.style={draw=cgray,fill=cgray!8,rounded corners=2pt,font=\sffamily\small,
    inner sep=4pt,align=center},
  pdec/.style={pbox,diamond,aspect=2.6,inner sep=0pt,draw=cgold,fill=cgold!12},
  parr/.style={-{Stealth[length=2mm]},thick,cgray}}

Bugs happen even with good tests. Debugging is the skill of finding out \emph{why} the program does what it does,
quickly and without guessing. This chapter goes from cheap tools (reading a traceback, \code{print}, f-strings) to
better ones (\code{assert}, \code{logging}) to the real debugger (\code{pdb}, \code{breakpoint()}).

% ------------------------------------------------------------
\section{Start here: read the traceback}

Debug from the bottom up first: you or a recent change broke it, so look at what you changed last.
If that fails, question what you trusted (a library, the data from an API); that is the slower, top-down search.

A \textbf{traceback} lists the chain of function calls that led to the error. Here is a program that fails when the API returns an
error payload instead of a normal reply (paths shortened with \code{...}):

\begin{pycode}
# usage_report.py
import json

def get_tokens(payload):
    return payload["usage"]["output_tokens"]

def report(raw):
    payload = json.loads(raw)
    return f"output tokens: {get_tokens(payload)}"

ok = '{"usage": {"output_tokens": 73}}'
error = '{"type": "error", "error": {"message": "overloaded"}}'

print(report(ok))
print(report(error))
\end{pycode}
\begin{rawcode}
output tokens: 73
Traceback (most recent call last):
  File ".../usage_report.py", line 17, in <module>
    print(report(error))
  File ".../usage_report.py", line 10, in report
    return f"output tokens: {get_tokens(payload)}"
  File ".../usage_report.py", line 5, in get_tokens
    return payload["usage"]["output_tokens"]
           ~~~~~~~^^^^^^^^^
KeyError: 'usage'
\end{rawcode}

(Real output. The \code{\textasciitilde} and \code{\textasciicircum} marker lines under the first two frames are left out.)

\begin{center}
%%MD-BEGIN
% ```text
% <module>      usage_report.py:17   print(report(error))        oldest call
%    | calls
%    v
% report()      usage_report.py:10   get_tokens(payload)
%    | calls
%    v
% get_tokens()  usage_report.py:5    payload["usage"][...]       <- failed here
%    v
% KeyError: 'usage'    <- WHAT went wrong (last line)
% ```
%%MD-END
\begin{tikzpicture}[node distance=5mm,fr/.style={pbox,text width=8.6cm,align=left,font=\ttfamily\footnotesize}]
  \node[fr] (a) {<module>\hfill usage\_report.py:17 \\ \textcolor{cgray}{print(report(error))}};
  \node[fr,below=of a] (b) {report()\hfill usage\_report.py:10 \\ \textcolor{cgray}{return f"output tokens: \{get\_tokens(payload)\}"}};
  \node[fr,below=of b,draw=cred,fill=cred!8] (c) {get\_tokens()\hfill usage\_report.py:5 \\ \textcolor{cgray}{return payload["usage"]["output\_tokens"]}};
  \node[pbox,below=of c,draw=cred,fill=cred!15,text width=8.6cm] (d) {\ttfamily\footnotesize\bfseries KeyError: 'usage'};
  \draw[parr] (a)--node[right,font=\sffamily\scriptsize]{calls}(b);
  \draw[parr] (b)--node[right,font=\sffamily\scriptsize]{calls}(c);
  \draw[parr,cred] (c)--(d);
  \node[font=\sffamily\scriptsize,text=cgray,anchor=west] at ($(a.east)+(3mm,0)$) {oldest call: where the program started};
  \node[font=\sffamily\scriptsize,text=cred,anchor=west] at ($(c.east)+(3mm,0)$) {the line that failed};
  \node[font=\sffamily\scriptsize,text=cred,anchor=west] at ($(d.east)+(3mm,0)$) {\textbf{what} went wrong};
\end{tikzpicture}
\end{center}

\begin{keyidea}[: how to read a traceback]
Read from the \textbf{bottom}. The last line names the error and its message. The frame just above it is the line that failed.
Walk upward until you find the first frame in \emph{your own} code; library frames are rarely the cause.
Here the real question is not ``why a \code{KeyError}'' but ``why did the API return no \code{usage}?'': the answer is an error payload.
\end{keyidea}

\begin{center}
%%MD-BEGIN
% ```mermaid
% flowchart TD
%     S["Something is wrong"] --> Q{"Exception with a traceback?"}
%     Q -- yes --> R["read the last line,<br/>then the last frame in YOUR code"]
%     Q -- no --> W["wrong output or hang:<br/>narrow it down with print / f-string / logging"]
%     R --> H{"cause clear?"}
%     W --> H
%     H -- no --> D["use the debugger:<br/>breakpoint() / pdb"]
%     H -- yes --> F["fix it"]
%     D --> F
%     F --> T["add a regression test"]
% ```
%%MD-END
\begin{tikzpicture}
  \node[pbox,draw=cred,fill=cred!10] (s) at (0.9,0) {something\\is wrong};
  \node[pdec] (q) at (3.9,0) {\footnotesize traceback?};
  \node[pbox,text width=3.1cm] (r) at (7.2,1.0) {read last line,\\then last frame\\in \emph{your} code};
  \node[pbox,text width=3.6cm] (w) at (7.4,-1.0) {wrong output or hang:\\print, f-string, logging};
  \node[pdec] (h) at (10.6,0) {\footnotesize cause clear?};
  \node[pbox] (d) at (10.6,-2.2) {debugger:\\\ttfamily breakpoint()};
  \node[pbox,draw=cgreen,fill=cgreen!10] (f) at (14.0,0) {fix it,\\add a test};
  \draw[parr] (s)--(q);
  \draw[parr] (q)--node[above,sloped,font=\scriptsize\sffamily]{yes}(r.west);
  \draw[parr] (q)--node[below,sloped,font=\scriptsize\sffamily]{no}(w.west);
  \draw[parr] (r.east)--(h.north west); \draw[parr] (w.east)--(h.south west);
  \draw[parr] (h)--node[above,font=\scriptsize\sffamily]{yes}(f);
  \draw[parr] (h)--node[right,font=\scriptsize\sffamily]{no}(d);
  \draw[parr] (d.east)-|(f.south);
\end{tikzpicture}
\end{center}

% ------------------------------------------------------------
\section{Assert}

\code{assert condition, message} states something that \emph{must} be true. If it is false, Python raises
\code{AssertionError} on the spot, with your message. Use it to catch impossible data early, before it causes a confusing
failure ten lines later.

\begin{pycode}
# assert_demo.py
def avg_tokens(counts):
    assert len(counts) > 0, "counts must not be empty"
    assert all(isinstance(c, int) for c in counts), f"ints only: {counts!r}"
    return sum(counts) / len(counts)

print(avg_tokens([10, 20]))
print(avg_tokens([]))
\end{pycode}
\begin{rawcode}
15.0
Traceback (most recent call last):
  ...
AssertionError: counts must not be empty
\end{rawcode}

\begin{trap}
Run Python with \code{-O} (\code{python -O assert\_demo.py}) and every \code{assert} is \emph{skipped}. The same call then fails
later with a less helpful error:
\begin{rawcode}
ZeroDivisionError: division by zero
\end{rawcode}
So use \code{assert} for internal sanity checks and tests. For input you do not control (user input, API data), check explicitly with
\code{if ...: raise ValueError(...)}.
\end{trap}

% ------------------------------------------------------------
\section{Print}

\code{print} is the fastest debugger. Print the values that matter, with labels. Three built-ins show variables:
\code{vars()} (or \code{locals()}) gives the local variables as a dict; \code{globals()} gives the module's.

\begin{pycode}
import sys

def func(*args, **kwargs):
    print(vars())
    print(sorted(locals()))

func(1, 2, model="small")
print("x" in globals())
print("debug info", file=sys.stderr)       # to the error stream
\end{pycode}
\begin{rawcode}
{'args': (1, 2), 'kwargs': {'model': 'small'}}
['args', 'kwargs']
True
debug info
\end{rawcode}

Debug output sent to \code{sys.stderr} stays separate from normal output: you can redirect your program's real output to a file and
still see the debug lines in the terminal.

\begin{trap}
Print debugging leaves clutter. Remove the prints before committing, or move them to \code{logging} (below), which you can switch off.
\end{trap}

% ------------------------------------------------------------
\section{F-strings}

An f-string inserts the value of any expression in \code{\{\}}. Adding \code{=} inside the braces prints the \emph{expression text and
its value}, so you never mislabel a debug line.

\begin{pycode}
x = "thing"
print(f"{x=}", f"{ x = }", f"{10/2=}")
tokens = 1234
print(f"{tokens=:,}")                   # with a format spec after the =
print(f"{type(tokens)=}")
\end{pycode}
\begin{rawcode}
x='thing'  x = 'thing' 10/2=5.0
tokens=1,234
type(tokens)=<class 'int'>
\end{rawcode}

Note \code{x='thing'}: the value is shown with \code{repr}, so strings have quotes, which makes stray spaces and \code{None} visible.

% ------------------------------------------------------------
\section{Pprint}

\code{pprint} (``pretty print'') lays out nested data over several lines, indented. Use it for JSON-like API replies.

\begin{pycode}
from pprint import pprint

text = "Hello there, how can I help you today?"
reply = {"type": "message", "role": "assistant",
         "content": [{"type": "text", "text": text}],
         "usage": {"input_tokens": 22, "output_tokens": 73}}
pprint(reply)
\end{pycode}
\begin{rawcode}
{'content': [{'text': 'Hello there, how can I help you today?',
              'type': 'text'}],
 'role': 'assistant',
 'type': 'message',
 'usage': {'input_tokens': 22, 'output_tokens': 73}}
\end{rawcode}

By default \code{pprint} sorts keys. Pass \code{sort\_dicts=False} to keep insertion order, and \code{width=60} to wrap sooner.
For real JSON you may prefer \code{json.dumps(obj, indent=2)}.

% ------------------------------------------------------------
\section{IceCream}

\textbf{IceCream} is a third-party package (\code{python -m pip install icecream}; not installed here, so the output below is
the documented style, not a captured run). Its \code{ic()} prints the expression, its value, and optionally the file, line
and function, and it \emph{returns} its argument, so you can wrap any sub-expression.

\begin{pycode}
from icecream import ic

tokens = 73
ic(tokens)                 # prints: ic| tokens: 73
ic(tokens * 2)             # prints: ic| tokens * 2: 146


def cost(n):
    ic()                   # prints file:line and the function name
    return n * 2


total = ic(cost(5)) + 1    # ic returns 10, so the expression still works
\end{pycode}

It is a nice shortcut, but \code{print(f"\{x=\}")} gives you most of the benefit with no install.

% ------------------------------------------------------------
\section{Decorators}

A decorator wraps a function to run code before or after it, without editing the function. That makes it a good place for debug output.
Here \code{trace} prints the arguments and the result of every call:

\begin{pycode}
import functools


def trace(func):
    """Print arguments and result of every call."""
    @functools.wraps(func)              # keep the original name and docstring
    def wrapper(*args, **kwargs):
        print(f"call {func.__name__} args={args} kwargs={kwargs}")
        result = func(*args, **kwargs)
        print(f"  -> {result!r}")
        return result
    return wrapper


@trace
def token_cost(model, tokens_in, tokens_out=0):
    rates = {"small": (0.10, 0.50)}
    r_in, r_out = rates[model]
    return (tokens_in * r_in + tokens_out * r_out) / 1e6


token_cost("small", 1000, tokens_out=500)
print(token_cost.__name__)
\end{pycode}
\begin{rawcode}
call token_cost args=('small', 1000) kwargs={'tokens_out': 500}
  -> 0.00035
token_cost
\end{rawcode}

\begin{typecheck}
After \code{@trace}, the name \code{token\_cost} refers to \code{wrapper}, a new function. \code{*args} collects positional arguments into a
tuple and \code{**kwargs} collects keyword arguments into a dict, so the wrapper accepts any call. Without \code{functools.wraps}
the name would print as \code{wrapper}.
\end{typecheck}

% ------------------------------------------------------------
\section{Logging}

\code{print} output vanishes when the terminal closes, and has no timestamp or importance. The standard \code{logging} module records
messages with a \textbf{level}, a time, and a name, and sends them where you choose. Use it in anything that runs longer than a quick script.

\subsection{Levels}

\begin{center}
%%MD-BEGIN
% ```text
% CRITICAL 50  -----------------------+
% ERROR    40                         |  shown when level=INFO
% WARNING  30  (default level)        |
% INFO     20  -----------------------+
% ---------------- level=INFO cuts here ----------------
% DEBUG    10  hidden
% ```
%%MD-END
\begin{tikzpicture}
  \foreach \i/\nm/\num/\c/\use in {0/CRITICAL/50/cred/{when the program cannot continue},
      1/ERROR/40/cred/{a call failed; the program goes on},
      2/WARNING/30/cgold/{something odd; the default threshold},
      3/INFO/20/cblue/{normal progress: ``got reply, 73 tokens''},
      4.5/DEBUG/10/cgray/{details: full request, variables}}{
    \node[pbox,draw=\c,fill=\c!10,font=\ttfamily\footnotesize,minimum width=2.4cm,anchor=west] at (0,-\i*0.52) {\nm\ \num};
    \node[font=\sffamily\footnotesize,anchor=west] at (2.7,-\i*0.52) {\use};}
  \draw[dashed,cgray] (0,-1.95)--(11.5,-1.95);
  \node[font=\sffamily\scriptsize,text=cgray,anchor=south east] at (11.5,-1.95) {\texttt{level=INFO} cuts here};
  \node[font=\sffamily\scriptsize,text=cgray,anchor=north east] at (11.5,-1.95) {below the line: hidden};
\end{tikzpicture}
\end{center}

Before any setup, only \code{WARNING} and above appear. This explains the book's puzzle of \code{debug} and \code{info} ``doing nothing'':

\begin{pycode}
import logging

logging.debug("not shown: default level is WARNING")
logging.info("not shown either")
logging.warning("shown")
logging.error("shown too")
\end{pycode}
\begin{rawcode}
WARNING:root:shown
ERROR:root:shown too
\end{rawcode}

(The functions \code{logging.warn} and \code{logging.fatal} are old, deprecated spellings; use \code{warning} and \code{critical}.)

\subsection{Configure once, with named loggers}

Call \code{basicConfig} \emph{once}, at the start of the program. Then each module asks for its own logger by name.
Names with dots form a family tree: \code{"chat.api"} is a child of \code{"chat"}, which is a child of the root logger.

\begin{pycode}
# log2.py
import logging

logging.basicConfig(level=logging.DEBUG,
                    format="%(levelname)-8s %(name)s:%(lineno)d %(message)s")
log = logging.getLogger("chat.api")

log.debug("sending request, model=%s", "small")
log.info("got reply, output_tokens=%d", 73)
log.warning("retrying after 429")
try:
    {}["usage"]
except KeyError:
    log.exception("could not read usage")     # message + the traceback
print("parent:", log.parent.name)
\end{pycode}
\begin{rawcode}
DEBUG    chat.api:7 sending request, model=small
INFO     chat.api:8 got reply, output_tokens=73
WARNING  chat.api:9 retrying after 429
ERROR    chat.api:13 could not read usage
Traceback (most recent call last):
  File ".../log2.py", line 11, in <module>
    {}["usage"]
    ~~^^^^^^^^^
KeyError: 'usage'
parent: root
\end{rawcode}

Two habits are visible here. Pass values as arguments (\code{"model=\%s", "small"}) instead of building the string yourself, so
Python formats it only if the message is actually emitted. Use \code{log.exception} inside \code{except}: it logs the traceback too.

\begin{trap}
\code{basicConfig} does nothing if the root logger already has handlers (an earlier call, or a library did it). Your level or format
silently stays as before. Pass \code{force=True} to reset. Real output of three calls:
\begin{pycode}
import logging
logging.basicConfig(level=logging.INFO, format="A %(message)s")
logging.basicConfig(level=logging.DEBUG, format="B %(message)s")   # ignored
logging.debug("debug?")                  # level is still INFO: hidden
logging.info("info")
logging.basicConfig(level=logging.DEBUG, format="C %(message)s", force=True)
logging.debug("debug now")
\end{pycode}
\begin{rawcode}
A info
C debug now
\end{rawcode}
\end{trap}

\subsection{Handlers and formatters}

Logging has four parts. A \textbf{logger} is what your code calls. A \textbf{handler} decides \emph{where} a record goes
(screen, file, email, web service). A \textbf{formatter} decides how it looks. A \textbf{filter} can drop records.

\begin{center}
%%MD-BEGIN
% ```mermaid
% flowchart LR
%     C["your code<br/>log.info(...)"] --> L["logger 'chat'<br/>level check"]
%     L --> H1["handler: console<br/>level INFO"] --> F1["formatter"] --> D1["terminal"]
%     L --> H2["handler: file<br/>level DEBUG"] --> F2["formatter"] --> D2["debug.log"]
% ```
%%MD-END
\begin{tikzpicture}[node distance=5mm]
  \node[pbox,draw=cblue,fill=cblue!10] (c) {your code\\{\ttfamily\footnotesize log.info(...)}};
  \node[pbox,right=of c] (l) {logger\\{\ttfamily\footnotesize "chat"}};
  \node[pbox,right=12mm of l,yshift=7mm] (h1) {handler: console\\{\footnotesize level INFO}};
  \node[pbox,right=12mm of l,yshift=-7mm] (h2) {handler: file\\{\footnotesize level DEBUG}};
  \node[pbox,right=of h1] (f1) {formatter};
  \node[pbox,right=of h2] (f2) {formatter};
  \node[pbox,right=of f1,draw=cgreen,fill=cgreen!10] (d1) {terminal};
  \node[pbox,right=of f2,draw=cgreen,fill=cgreen!10] (d2) {\ttfamily\footnotesize debug.log};
  \draw[parr] (c)--(l); \draw[parr] (l)--(h1); \draw[parr] (l)--(h2);
  \draw[parr] (h1)--(f1); \draw[parr] (h2)--(f2); \draw[parr] (f1)--(d1); \draw[parr] (f2)--(d2);
\end{tikzpicture}
\end{center}

\begin{pycode}
# log4.py
import logging
import sys

log = logging.getLogger("chat")
log.setLevel(logging.DEBUG)                  # logger lets everything through

console = logging.StreamHandler(sys.stdout)
console.setLevel(logging.INFO)               # but the screen shows INFO+
console.setFormatter(logging.Formatter("%(levelname)s %(message)s"))

logfile = logging.FileHandler("debug.log", encoding="utf-8")
logfile.setLevel(logging.DEBUG)              # the file keeps everything
logfile.setFormatter(logging.Formatter("%(levelname)s %(name)s %(message)s"))

log.addHandler(console)
log.addHandler(logfile)

log.debug("prompt has 12 messages")
log.info("reply received")
log.error("rate limited")
\end{pycode}
\begin{rawcode}
INFO reply received
ERROR rate limited
\end{rawcode}

The file \code{debug.log} then contains all three:

\begin{rawcode}
DEBUG chat prompt has 12 messages
INFO chat reply received
ERROR chat rate limited
\end{rawcode}

For a simple one-file setup, \code{basicConfig(filename="app.log", level=..., format=...)} is enough. A format such as
\code{"\%(asctime)s \%(levelname)s \%(message)s"} adds a time stamp, which looks like \code{2026-11-02 13:02:01,655 ERROR boom}.
Other useful fields: \code{\%(name)s}, \code{\%(lineno)d}, \code{\%(funcName)s}.

\begin{gamerule}[: logging in an LLM app]
Log the model name, token counts, latency and the request id, at \code{INFO}. Log the full prompt only at \code{DEBUG}, because
prompts can contain private data. \textbf{Never log an API key.} Keep the file out of Git (\code{.gitignore}).
\end{gamerule}

% ------------------------------------------------------------
\section{Pdb: the debugger}

A debugger runs your program under your control: you pause it, look at variables, and step through lines.
Python's built-in debugger is \code{pdb}. Here is a data-dependent bug. The function should total tokens per model from lines like
\code{"small,100"}, but the total for \code{small} comes out wrong.

\begin{pycode}
# totals.py
def total_usage(lines):
    totals = {}
    for line in lines:
        model, tokens = line.split(",")
        totals[model] = int(tokens)
    return totals


if __name__ == "__main__":
    data = ["small,100", "large,40", "small,50"]
    print(total_usage(data))
\end{pycode}
\begin{rawcode}
{'small': 50, 'large': 40}
\end{rawcode}

No exception, so there is no traceback to read: it is a \emph{logic} error. \code{small} should be 150. Start the program under pdb
(\code{-m pdb}), set a breakpoint on line 5 with \code{b 5}, and continue with \code{c}:

\begin{rawcode}
> python -m pdb totals.py
> ...\totals.py(1)<module>()
-> def total_usage(lines):
(Pdb) b 5
Breakpoint 1 at ...\totals.py:5
(Pdb) c
> ...\totals.py(5)total_usage()
-> totals[model] = int(tokens)
(Pdb) c
> ...\totals.py(5)total_usage()
-> totals[model] = int(tokens)
(Pdb) c
> ...\totals.py(5)total_usage()
-> totals[model] = int(tokens)
(Pdb) p line
'small,50'
(Pdb) p totals
{'small': 100, 'large': 40}
(Pdb) n
> ...\totals.py(3)total_usage()
-> for line in lines:
(Pdb) p totals
{'small': 50, 'large': 40}
\end{rawcode}

(Real session; only the long path is shortened.) After the third stop, \code{totals} holds \code{small: 100}. The next \emph{single} line
replaces it with 50: the code \emph{assigns} (\code{=}) where it should \emph{add}. Fix it, and turn the bug into a regression test:

\begin{pycode}
from collections import defaultdict


def total_usage(lines):
    totals = defaultdict(int)
    for line in lines:
        model, tokens = line.split(",")
        totals[model] += int(tokens)        # add, do not overwrite
    return dict(totals)


def test_repeated_model_is_added():
    data = ["small,100", "large,40", "small,50"]
    assert total_usage(data) == {"small": 150, "large": 40}


test_repeated_model_is_added()
print(total_usage(["small,100", "large,40", "small,50"]))
\end{pycode}
\begin{rawcode}
{'small': 150, 'large': 40}
\end{rawcode}

\subsection{Commands you will use}

\begin{center}
\begin{tabularx}{\linewidth}{@{}l l X l@{}}
\toprule
\textbf{pdb} & \textbf{Name} & \textbf{Does} & \textbf{VS Code} \\
\midrule
\code{n} & next & run this line, \emph{step over} any call & F10 \\
\code{s} & step & \emph{step into} the call on this line & F11 \\
\code{r} & return & run until the current function returns & Shift+F11 \\
\code{c} & continue & run until the next breakpoint or the end & F5 \\
\code{b 5} / \code{b} & breakpoint & set one on line 5 / list them all & click the gutter \\
\code{l} / \code{ll} & list & show code around the line / the whole function & the editor \\
\code{p expr} / \code{pp expr} & print & show a value / pretty-print it & Variables panel \\
\code{w}, \code{u}, \code{d} & where, up, down & show the call stack, move between frames & Call Stack panel \\
\code{q} & quit & stop the program and leave & Shift+F5 \\
\bottomrule
\end{tabularx}
\end{center}

\begin{center}
%%MD-BEGIN
% ```text
%  line 8:  total = add(a, b)      <- paused here
%                |
%     n ---------+--------> line 9   (add() ran in one go)
%     s ---------+--------> line 2 inside add()   (you are now in add)
% ```
%%MD-END
\begin{tikzpicture}[node distance=7mm]
  \node[pbox,draw=cgold,fill=cgold!12,font=\ttfamily\footnotesize,text width=4.3cm,align=left] (a) {8: total = add(a, b)};
  \node[pbox,right=24mm of a,yshift=7mm,font=\ttfamily\footnotesize,text width=4.2cm,align=left] (b) {9: print(total)};
  \node[pbox,right=24mm of a,yshift=-7mm,font=\ttfamily\footnotesize,text width=4.2cm,align=left] (c) {2: return a + b\\{\sffamily\scriptsize inside add()}};
  \draw[parr] (a.east)--node[above,sloped,font=\ttfamily\scriptsize]{n: step over}(b.west);
  \draw[parr] (a.east)--node[below,sloped,font=\ttfamily\scriptsize]{s: step into}(c.west);
  \node[font=\sffamily\scriptsize,text=cgray,above=1mm of a] {paused here (the arrow \texttt{->} in pdb)};
\end{tikzpicture}
\end{center}

\begin{gamerule}
Use \code{n} when you trust a function (library code, code you tested) and \code{s} when you suspect it. Do not step line by line from the
top: set a breakpoint near the suspect and \code{c} to it.
\end{gamerule}

\subsection{Post-mortem: debug a crash where it happened}

If a program crashes, run it under \code{python -m pdb script.py} and type \code{c}. On the uncaught exception pdb stops at the failing frame (``post mortem''),
with all variables intact. For \code{usage\_report.py} that looks like this (condensed from a real session):

\begin{rawcode}
(Pdb) c
Uncaught exception. Entering post mortem debugging
> ...\usage_report.py(5)get_tokens()
-> return payload["usage"]["output_tokens"]
(Pdb) p payload
{'type': 'error', 'error': {'message': 'overloaded'}}
(Pdb) up
> ...\usage_report.py(10)report()
(Pdb) p raw
'{"type": "error", "error": {"message": "overloaded"}}'
\end{rawcode}

Now you can \emph{see} that the payload is an error object. The fix is to check before reading it:

\begin{pycode}
def get_tokens(payload):
    if payload.get("type") == "error":
        raise RuntimeError(payload["error"]["message"])
    return payload["usage"]["output_tokens"]
\end{pycode}

Another useful flag is \code{python -i script.py}: after the script ends or crashes, you stay in the interactive prompt and can inspect the variables.

% ------------------------------------------------------------
\section{\texttt{breakpoint()}}

Add \code{breakpoint()} to your code and Python opens pdb \emph{at that line}, with no setup. It is the quickest way into the debugger.

\begin{pycode}
# bp.py
def total_usage(lines):
    totals = {}
    for line in lines:
        model, tokens = line.split(",")
        breakpoint()                       # pause here on every loop
        totals[model] = totals.get(model, 0) + int(tokens)
    return totals

print(total_usage(["small,100", "small,50"]))
\end{pycode}
\begin{rawcode}
> ...\bp.py(6)total_usage()
-> totals[model] = totals.get(model, 0) + int(tokens)
(Pdb) p model, tokens
('small', '100')
(Pdb) c
> ...\bp.py(6)total_usage()
-> totals[model] = totals.get(model, 0) + int(tokens)
(Pdb) p totals
{'small': 100}
(Pdb) c
{'small': 150}
\end{rawcode}

The environment variable \code{PYTHONBREAKPOINT} chooses what \code{breakpoint()} does. Setting it to \code{0} turns every breakpoint off
without editing the code (real output of the same run):

\begin{rawcode}
> set PYTHONBREAKPOINT=0                (cmd.exe)
> $env:PYTHONBREAKPOINT = "0"           (PowerShell)
> python bp.py
{'small': 150}
\end{rawcode}

Setting it to the name of another debugger function (for example \code{web\_pdb.set\_trace}, from a package you install) switches debugger.

\begin{trap}
Never commit a \code{breakpoint()} call: a server or a CI job would stop and wait for input forever. Search for it before every commit
(\code{git grep breakpoint}), or let your linter flag it.
\end{trap}

\begin{boss}[: find the cause, do not guess]
Programmers lose hours changing code at random. A better loop has four steps:
(1) \textbf{reproduce} the bug with the smallest input that still fails (three lines of data, not the whole file);
(2) \textbf{state a hypothesis} about the cause (``the dict value is overwritten, not added'');
(3) \textbf{test the hypothesis} by observing, not changing: \code{print}, a log line, or \code{p totals} in pdb at the right moment;
(4) \textbf{fix one thing, then add a test} that fails without the fix. If the observation does not match your hypothesis, your mental model is wrong;
revise it instead of editing more code. Most hard bugs are a wrong assumption about the data, not a wrong line of code.
\end{boss}

\begin{practice}
\begin{enumerate}[leftmargin=*,itemsep=2pt]
\item A traceback ends with \code{KeyError: 'usage'}. Which line of the traceback do you read first, and what two questions do you ask next?
\item Use an f-string with \code{=} to print the expression \code{len(messages) * 2} and its value, and the type of the variable \code{reply}.
\item Why does \code{logging.debug("x")} print nothing in a fresh script? Give two different ways to make it appear.
\item Set up logging so that INFO and above go to the screen but DEBUG and above also go to a file. Which objects do you create?
\item You run \code{python -O} and a check vanishes. Which statement was it, and what should you use for checking user input instead?
\item In pdb, you are on a line that calls \code{summarize(client, text)}. What is the difference between typing \code{n} and \code{s}?
\end{enumerate}
\end{practice}

\begin{achievement}{Bug Hunter}
\begin{itemize}[leftmargin=*,itemsep=1pt]
\item Read a traceback from the bottom and locate the failing line in your own code.
\item Debug with f-string \code{=}, \code{pprint}, \code{assert} and a tracing decorator.
\item Configure \code{logging} with levels, named loggers, a file handler and a format.
\item Step through code in \code{pdb}, set breakpoints, inspect variables, and debug a crash post-mortem.
\item Turn a fixed bug into a regression test.
\end{itemize}
\end{achievement}

%%MD-BEGIN
% **Related:** [[10 - Functions]] (decorators, `*args`), [[04 - Strings]] (f-strings), [[09 - Dictionaries and Sets]] (`KeyError`), [[13 - Development Environment]] (which Python?), [[14 - Type Hints and Documentation]] (catch bugs earlier), [[15 - Testing]] (regression tests), [[20 - Files]] (log files).
%%MD-END
